\begin{table}[H]
  \centering
  \caption{Trayectorias a 3 olas (2010$\to$2013$\to$2016) de hogares
  NO pobres en 2010, panel emparejado 1 a 1, monetaria vs. IPM.}
  \label{tab:trayectorias_3olas}
  \footnotesize
  \begin{tabular}{lcccc}
    \toprule
    & \multicolumn{2}{c}{\textbf{Monetaria}} & \multicolumn{2}{c}{\textbf{IPM}} \\
    \textbf{Trayectoria} & \textbf{n} & \textbf{\%} & \textbf{n} & \textbf{\%} \\
    \midrule
    Nunca cae & 1{,}700 & 67.2\% & 4{,}381 & 81.2\% \\
    Cae tarde (solo en 2016) & 220 & 8.7\% & 357 & 6.6\% \\
    Entra y sale (transitorio) & 355 & 14.0\% & 445 & 8.2\% \\
    Entra y se queda (persistente) & 255 & 10.1\% & 214 & 4.0\% \\
    \bottomrule
  \end{tabular}
  \begin{minipage}{0.85\textwidth}
    \vspace{4pt}
    \footnotesize \textit{Nota:} $n$=2{,}530 hogares no pobres en 2010 bajo
    pobreza monetaria; $n$=5{,}397 bajo IPM (universos distintos porque
    no-pobre-monetaria y no-pobre-IPM en 2010 no son el mismo conjunto de
    hogares). \% sobre ese universo, sin ponderar (ver también
    \texttt{pct\_ponderado} en los CSV fuente). Fuente: cálculos propios,
    \texttt{outputs/tables/eda\_transicion\_covariables/trayectorias\_3olas\_\{monetaria,ipm\}.csv}
    (\texttt{src/02\_build/eda\_trayectorias\_3olas.py}), generados por
    \texttt{src/tabla\_trayectorias\_3olas.py}.
  \end{minipage}
\end{table}